% !TEX root = ../main.tex
\section{Bregman divergences for value learning}
\label{sec:04_bregman}

% ---------------------------------------------------------------------------
\subsection{Any Bregman divergence learns the value function}

For a strictly convex, differentiable potential $\varphi$, the Bregman
divergence is
\begin{equation}
  \Breg(u, v) = \varphi(u) - \varphi(v) - \varphi'(v)\,(u - v).
\end{equation}
Squared error is the special case $\varphi(u)=u^2$. The property we exploit:

\begin{lemma}[Bregman divergences are proper losses for the value]
\label{lem:minimizer}
Let the regression target be $e^{q}$ with $q$ a conditionally-unbiased
estimate of the value (the H-martingale target, \S\ref{sec:02_setting}). For
any strictly convex $\varphi$,
\begin{equation}
  \arg\min_{u}\ \E_{q\mid x_t}\!\big[\Breg(u,\, e^{q})\big] \;=\; \E[e^{q}\mid x_t],
\end{equation}
so the minimizing prediction in log space is $p^\star=\log\E[e^{q}\mid x_t]=\Vf(x_t)$.
\end{lemma}
\noindent Thus the choice of $\varphi$ does \emph{not} bias the learned value;
it is free to be chosen for \emph{conditioning}. (Proof in
App.~\ref{sec:D_proofs}.)

% ---------------------------------------------------------------------------
\subsection{Design criteria and two reference points}

We parameterize the prediction in log space, $u=e^{p}$, and ask how the
gradient $\partial L/\partial p$ behaves as the log-error $p-q\to\pm\infty$.

\paragraph{Squared error ($\varphi(u)=u^2$).}
$\partial L/\partial p = 2e^{p}(e^{p}-e^{q}) = 2e^{2p}-2e^{p+q}$, which
\emph{explodes} as $p\to+\infty$ and \emph{vanishes} as $p\to-\infty$ --- the
under-prediction pathology of \S\ref{sec:03_expmse}.

\paragraph{Itakura--Saito ($\varphi(u)=-\ln u$).}
$D_{\mathrm{IS}}(e^{q},e^{p}) = e^{q-p}-(q-p)-1$; its gradient in $p$ is
$1-e^{q-p}$, a \emph{linear} (bounded-gradient) tail for over-predictions that
removes the vanishing problem, but it still explodes for under-predictions.
It is \emph{scale-invariant} ($D_{\mathrm{IS}}(cu,cv)=D_{\mathrm{IS}}(u,v)$)
and is the negative log-likelihood of a gamma/exponential observation model,
giving a natural uncertainty interpretation for $\Vf$. A useful midpoint, but
one-sided.

% ---------------------------------------------------------------------------
\subsection{The Spence loss}
\label{subsec:spence}

We propose the Bregman divergence generated by the potential $F$ with
\begin{equation}
  F''(u) = -\,\frac{\ln u}{u\,(1-u)},
\end{equation}
applied to $u=e^{p}$ against the target $e^{q}$. Its gradient with respect to
the prediction has the closed form
\begin{equation}
  \boxed{\ \frac{\partial L}{\partial p} \;=\; p\,\frac{e^{p}-e^{q}}{e^{p}-1}\ }
  \label{eq:spence-grad}
\end{equation}
which is zero iff $p=q$ and, by Lemma~\ref{lem:minimizer}, has population
minimizer $\Vf$.

\paragraph{Two-sided quadratic tails.}
From \eqref{eq:spence-grad},
\begin{align}
  p\to+\infty:\quad & \frac{\partial L}{\partial p}\to p
    &&\text{(no $e^{2p}$ blow-up; over-prediction is quadratically penalized),}\\
  p\to-\infty:\quad & \frac{\partial L}{\partial p}\to p\,e^{q}
    &&\text{(non-vanishing corrective signal for under-prediction).}
\end{align}
Both tails are linear in $p$, i.e.\ the loss is quadratic on both sides --- the
first loss to fix both pathologies of exponential-space regression at once.

\paragraph{Value via the Spence function.}
The divergence value (used only for monitoring; the gradient uses
\eqref{eq:spence-grad}) is expressible through the Spence function
$S(x)=\Spence(1-e^{x})$, using the dilogarithm identity
$S'(x)=-\,x e^{x}/(e^{x}-1)=x/\mathrm{expm1}(-x)$; hence the name. Full
derivation, the numerically stable branch-wise gradient (finite for all float
inputs, unlike the naive autograd which NaNs for $p\gtrsim 88$), and the
gamma-like uncertainty interpretation are in App.~\ref{sec:A_spence}.

% ---------------------------------------------------------------------------
\begin{table}[t]
\centering
\caption{The value-learning loss family. Tails give $\partial L/\partial p$ as
the log-error $\to\pm\infty$. All three share the same minimizer,
$\Vf=\log\E[e^{q}\mid x_t]$ (Lemma~\ref{lem:minimizer}).}
\label{tab:loss-family}
\begin{tabular}{lcccc}
\toprule
loss & under-predict ($p\!\to\!-\infty$) & over-predict ($p\!\to\!+\infty$)
     & scale-inv. & uncertainty model \\
\midrule
squared error   & vanishes ($e^{p+q}$)     & explodes ($e^{2p}$) & no  & Gaussian (exp) \\
Itakura--Saito  & explodes                 & linear (bounded)    & yes & gamma \\
\textbf{Spence (ours)} & \textbf{quadratic} ($p\,e^{q}$) & \textbf{quadratic} ($p$) & no & gamma-like \\
\bottomrule
\end{tabular}
\end{table}
